\documentclass[11pt]{article}

\usepackage[a4paper,margin=2.5cm]{geometry}
\usepackage{fontspec}
\usepackage{unicode-math}
\setmainfont{lmroman10-regular.otf}[BoldFont=lmroman10-bold.otf,ItalicFont=lmroman10-italic.otf,BoldItalicFont=lmroman10-bolditalic.otf]
\setmathfont{latinmodern-math.otf}
\setmonofont{Menlo}[Scale=0.82]
\usepackage{microtype}
\usepackage{amsthm}
\usepackage{mathtools}
\usepackage{booktabs}
\usepackage{tabularx}
\usepackage{longtable}
\usepackage{enumitem}
\usepackage{fancyvrb}
\usepackage{xcolor}
\usepackage[numbers,sort&compress]{natbib}
\usepackage{xurl}
\usepackage[colorlinks=true,linkcolor=blue!50!black,citecolor=green!35!black,urlcolor=blue!60!black]{hyperref}
\usepackage{caption}
\hypersetup{pdftitle={A Candidate Lean 4 Statement of the Yang-Mills Existence and Mass Gap Problem},pdfauthor={Peter Lowes}}

\newtheorem{proposition}{Proposition}[section]
\theoremstyle{definition}
\newtheorem{definition}[proposition]{Definition}
\theoremstyle{remark}
\newtheorem{remark}[proposition]{Remark}

\fvset{fontsize=\small,frame=single,framesep=2mm}
\newcommand{\R}{\mathbb{R}}
\newcommand{\C}{\mathbb{C}}
\newcommand{\Sch}{\mathcal{S}}
\newcommand{\tr}{\operatorname{tr}}
\newcommand{\lean}[1]{\texttt{\hyphenchar\font=-1\relax #1}}
\setlength{\emergencystretch}{4em}
\newcommand{\dotbr}{.\allowbreak}

\title{A Candidate Lean~4 Statement of the Yang--Mills\\ Existence and Mass Gap Problem\\[4pt]
\large Adversarial AI Review and a Research Programme\\ (No Proof Is Claimed)}
\author{Peter Lowes}
\date{6 November 2026; revised 10--11 November 2026 \quad (version 0.2.0)\\[2pt]\small Statement files unchanged since 6 November; pre-registration claim withdrawn after\\ the release of OpenAI's Lean code on 10 November.}

\begin{document}
\maketitle

\begin{abstract}
On 26 October 2026 OpenAI announced an AI-generated proof of the existence of four-dimensional
quantum Yang--Mills theory, formalized in Lean, and placed it with a shepherding committee whose
first stated objective is to check that ``the definitions and problem statement correspond with
those which are standard in physics and mathematics''. On 10 November 2026 the committee published OpenAI's Lean code, with an assessment that
its announcement describes as a validation that the code proves the existence of pure Yang--Mills
theory~\cite{Douglas2026g}. That release came after the statement, its review and the research
library reported here were completed on 6 November; this version examines neither the code nor the
assessment. This paper does \emph{not} prove existence or the
mass gap. It contributes four things that are independent of OpenAI's material.
(i)~A candidate formal statement, in Lean~4 with Mathlib, of the Clay Millennium problem
``Yang--Mills existence and mass gap'' and of an existence-only variant. The statement uses the
Osterwalder--Schrader axioms with OS~II's linear growth condition, a Euclidean formulation of
the mass gap, and identifies the theory as Yang--Mills through a continuum-limit clause for
Wilson's lattice gauge theory, a full-field-content clause for the 21 components of
$\tr F_{\mu\nu}F_{\kappa\lambda}$, and a short-distance asymptotic-freedom fingerprint.
(ii)~Machine-checked sanity and non-vacuity results: the trivial theory satisfies the axioms but
not the non-triviality clause; the lattice partition function is positive; and $SU(N)$,
$N\ge2$, satisfies the hypotheses of the general statement ($SU(N)$ is connected and its Lie
algebra, computed in Lean to be $\mathfrak{su}(N)$, is simple). All of these depend only on
Lean's three standard axioms.
(iii)~An adversarial-review methodology, executed with AI agents, that found a critical defect in
the first draft (the cluster axiom admitted the zero vector, which collapses the Hilbert space
and would have made the open theorems provably false) and four major fidelity defects, two
fixed and two only partly (short-distance agreement is imposed only through a scalar two-point
fingerprint, and the dependence on Wilson's lattice regularization remains); both residues are
recorded as deviations from the Clay text.
(iv)~A research library compiled and cross-checked by AI agents (11 notes; 666 bibliography
entries whose existence and metadata were checked, a few flagged as unconfirmed) and a research programme
for an independent proof, informed by the published record of OpenAI's agent-swarm attack on
Navier--Stokes and its public Lean repositories, with explicit subjective odds, stop rules and
an independence policy. The statement's main design elements also appear in an earlier candidate
statement by R\'emondi\`ere~\cite{Remondiere2026}. AI systems wrote all of the mathematics, Lean
code and prose under the author's direction; no human expert has reviewed the work;
Section~\ref{sec:ai} gives a full disclosure.
\end{abstract}

\tableofcontents

%=====================================================================
\section{Introduction}
%=====================================================================

The Clay Mathematics Institute's Millennium problem on Yang--Mills theory, described by Jaffe
and Witten~\cite{Jaffe2000}, asks for a proof that for every compact simple gauge group $G$ a
non-trivial quantum Yang--Mills theory exists on $\R^4$ and has a mass gap $\Delta>0$, where
existence includes ``axiomatic properties at least as strong as'' those of
Streater--Wightman~\cite{Streater1964} and Osterwalder--Schrader~\cite{Osterwalder1973,Osterwalder1975}.
Douglas's 2004 status report~\cite{Douglas2004} recorded the state of the art as
Ba{\l}aban's renormalization-group programme for lattice gauge theory~\cite{Balaban1984a,Balaban1987,Balaban1988d,Balaban1989b,Balaban1989c}
and the infrared-cutoff construction of Magnen, Rivasseau and S\'en\'eor~\cite{Magnen1993}; a
2026 review is~\cite{Douglas2026f}.

\paragraph{The 2026 situation.}
On 8 September 2026 OpenAI announced a Lean-formalized proof of finite-time blowup for the
three-dimensional Navier--Stokes equations with smooth forcing, produced by a large multi-agent
system~\cite{OpenAI2026g,OpenAI2026c,OpenAI2026k,Howlett2026h}; the scope of that result relative to
Fefferman's problem statement~\cite{Fefferman2000} is disputed because it uses the forced
alternative~\cite{Howlett2026b}; Constantin, Ignatova and Vicol show that constructions with its key
features cannot have a real-analytic force~\cite{Constantin2026c}.
On 26 October 2026 OpenAI announced, among other results, an existence proof for pure
Yang--Mills theory in four dimensions, as part of a claimed construction of certain interacting
gauge theories with fermions; the mass gap is not claimed and no gauge group is named
publicly~\cite{OpenAI2026p,Woit2026d,QRC2026,Howlett2026e}. The natural-language proofs (``the full
1147-page LLM proofs'' according to the committee, covering an unstated subset of the claimed results)
and the Lean formalization of pure Yang--Mills existence are held by a shepherding committee (Douglas, Gross (chair), Jaffe, Klebanov, Nekrasov), which states that
Ba{\l}aban's ideas were ``among the important elements'' and that confinement and the mass gap
remain open~\cite{Douglas2026e,Douglas2026d,KITP2026b}. As of 6 November 2026, when the statement, its review and
the research library reported here were completed, none of this material was public. On 10 November
2026 the committee published the Lean code that OpenAI submitted to it on 28 October, together with
an assessment~\cite{Douglas2026g}. The committee's announcement calls the assessment a ``Validation of
the OpenAI Yang--Mills construction (Lean proof)'' and describes it as summarizing the committee's
checks that the code proves the existence of pure Yang--Mills theory according to the problem
statements standard in physics and mathematics; that is the committee's conclusion, not one examined
here. The announcement covers only the Lean code and the assessment. GitHub lists the repository
holding OpenAI's code as created on 5 November 2026; the committee announced that it would make it
public on 10 November. Neither the code nor the assessment is examined in this version. Independently, Sussman and Alp\"oge have posted AI-generated, self-described tentative
manuscripts continuing Ba{\l}aban's programme in three dimensions, with an unfinished
four-dimensional part~\cite{Sussman2026}, and Aronow and Lopatto have posted AI-generated
manuscripts claiming the continuum limit of three-dimensional $SU(2)$ Yang--Mills theory on a
torus~\cite{Aronow2026}. None of these claims has
been refereed.

\paragraph{Why a formal statement matters now.}
A Lean proof certifies only the theorem that was stated, not the paper's argument: the
Navier--Stokes Lean development reaches its headline theorem through intermediate lemmas that differ
from the paper's~\cite{Bastounis2026}. Errors of meaning therefore concentrate in the statement. The
maintainers of the lean-dojo collection mark its Yang--Mills statement as incomplete and ``not a valid
target'', noting that a two-dimensional toy model satisfies it~\cite{Leandojo2026}. An earlier
AI-written candidate statement by R\'emondi\`ere, deposited on Zenodo on 25 October
2026~\cite{Remondiere2026}, already combines a Wilson-lattice-limit variant with an axiomatic variant
carrying a Yang--Mills ``fingerprint'', uses the 21 components of $\tr F_{\mu\nu}F_{\rho\sigma}$ as
field content, and concedes that its gap covers only the C-even sector they generate; like the present
work, it was written by Claude agents under its author's direction. Our statement shares these main
design elements. The library note that reviewed it (note~01) was written before our version~0.1 was
reviewed, and our version~0.2 additions were made by reviewer agents who had read that note. As of
6 November 2026 the formal-conjectures collection~\cite{DeepMind2026} contains no merged Yang--Mills
statement (an open pull request proposes one). The shepherding committee lists statement fidelity as
its first objective~\cite{Douglas2026e}, and its 9 November announcement reports a positive result of
checks of that kind~\cite{Douglas2026g}. An AI-reviewed candidate statement written independently of
OpenAI's code, with published file hashes, is useful as a point of comparison whatever the status of
that proof. It is \emph{not} a pre-registration. The statement files of version~0.2.0 were completed
on 6 November 2026, four days before the committee published OpenAI's Lean code, but that date rests
on the author's account and on local Git commits that carry no independent timestamp. Those commits
were squashed into a single commit before first publication, to remove local file paths, so they are
not public: the published history begins on 10 November 2026, and this deposit was made after the
release. The author states that the statement files were not changed after the release and that
OpenAI's Lean code was not consulted in writing them; neither claim can be checked independently. What the deposit does fix is the content: the
SHA-256 hashes of the statement files are listed in Section~\ref{sec:data}, and any later version will
be published with a diff against 0.2.0 and the reason for each change.

\paragraph{Premise and scope.}
This work was undertaken under a working premise set by the author: \emph{assume that OpenAI's
existence proof and its Lean formalization are correct, and aim to produce an independent proof
and formalization}. The premise is not evaluated here. Section~\ref{sec:programme} explains what
it does and does not imply for our formal targets.

\paragraph{Contributions.}
\begin{enumerate}[leftmargin=*]
\item A candidate Lean~4/Mathlib statement (Section~\ref{sec:statement}) of
  \lean{ClayYangMills}, its existence half \lean{ClayExistence}, an $SU(N)$ existence statement
  along the asymptotic-freedom trajectory (\lean{SUNExists}), and an ordered ladder of weaker
  milestone statements (implied by \lean{SUNExists}, with the implications proved in Lean). Its main
  design elements also appear in the earlier candidate of~\cite{Remondiere2026}.
\item Fully machine-checked sanity and non-vacuity results (Section~\ref{sec:checked}).
\item An adversarial statement-review methodology and its findings (Section~\ref{sec:review}).
\item An AI-compiled, AI-cross-checked research library and an explicit research programme with odds, stop rules and
  an AI operating model informed by the published record of 2026's large-scale AI attacks on
  Millennium problems (Sections~\ref{sec:landscape} and~\ref{sec:programme}).
\end{enumerate}
\emph{This paper contains no proof of existence or of a mass gap, and does not claim one.}

%=====================================================================
\section{The problem}\label{sec:problem}
%=====================================================================

\subsection{What Jaffe and Witten ask for}
The official description~\cite{Jaffe2000} asks for a quantum field theory ``in the above sense''
(Wightman/OS axioms~\cite{Streater1964,Osterwalder1973,Osterwalder1975}; see
also~\cite{Glimm1987,Haag1996}) satisfying, in addition:
\begin{enumerate}[leftmargin=*]
\item \textbf{Field content.} Local quantum fields correspond to gauge-invariant local polynomials
  in the curvature $F$ and its covariant derivatives, e.g.\ $\tr F_{ij}F_{kl}(x)$; the
  correspondence is not canonical (a footnote notes renormalization subtleties).
\item \textbf{Short distances.} Correlation functions agree at short distances with asymptotic
  freedom~\cite{Gross1973,Politzer1973} and textbook renormalized perturbation theory, including
  a stress tensor and an operator product expansion.
\item \textbf{Mass gap.} $H\Omega=0$, the spectrum of $H$ lies in $[0,\infty)$, and $H$ has no
  spectrum in $(0,\Delta)$; the supremum $m$ of such $\Delta$ must be finite.
\item \textbf{No weak existence.} Existence by compactness alone is not a solution unless the
  axioms and the gap are established for the limit.
\item \textbf{A Hilbert space.} ``Establishing a quantum mechanical Hilbert space is part of the
  solution.''
\end{enumerate}

\subsection{Osterwalder--Schrader axioms and linear growth}
Osterwalder and Schrader~\cite{Osterwalder1973} gave Euclidean axioms (E0)--(E4) for Schwinger
distributions $S_n$ on $\Sch_0(\R^{4n})$ (Schwartz functions vanishing with all derivatives on
all diagonals): temperedness, Euclidean covariance, reflection positivity, symmetry, and cluster
property. A lemma in~\cite{Osterwalder1973} is incorrect, and OS~II~\cite{Osterwalder1975}
replaced temperedness by the \emph{linear growth condition}
\begin{equation}\label{eq:E0prime}
\tag{E0$'$}
|S_n(f)|\le \sigma_n\,|f|_{ns},\qquad \sigma_n\le\alpha\,(n!)^\beta,
\end{equation}
for a fixed Schwartz norm index $s$, under which reconstruction of a Wightman theory holds. The
Glimm--Jaffe measure axioms~\cite{Glimm1987} are not suitable here. A short argument of ours (library
note~01, \S2.6; not a published result) shows that their regularity axiom OS1 is incompatible with any
field of scaling dimension greater than two, and so with the expected $|x|^{-8}\log^{-2}(1/|x|)$
behaviour of $\tr F^2$.

\subsection{Status}
Table~\ref{tab:status} summarises the status on 6 November 2026; details and sources are in the
library notes (Section~\ref{sec:data}).

\begin{table}[ht]
\centering\small
\begin{tabularx}{\textwidth}{@{}lX@{}}
\toprule
Item & Status \\
\midrule
4D existence (gauge groups not stated) & \textbf{Claimed} by OpenAI for pure Yang--Mills theory; gauge groups, axioms and Lean statement not public on 6 November; Lean code released by the committee on 10 November, after this work, with an assessment that the committee's announcement describes as validating it; neither is examined here~\cite{OpenAI2026p,Douglas2026e,Douglas2026d,Douglas2026g}. \\
4D mass gap & \textbf{Open}: not claimed by OpenAI~\cite{OpenAI2026p,Woit2026d}; the committee states that it remains open~\cite{Douglas2026e}. \\
3D Yang--Mills continuum & No refereed construction; AI-generated claims~\cite{Sussman2026,Aronow2026}. \\
2D Yang--Mills & Constructed~\cite{Driver1989,Gross1989,Sengupta1997,Levy2003,Chevyrev2019}; Langevin dynamics and universality~\cite{Chandra2022,Chevyrev2026b}. \\
4D lattice mass gap, continuous non-Abelian $G$ & Proved only at strong coupling~\cite{Osterwalder1978,Seiler1982,Shen2023}; never in the weak-coupling continuum regime. (Finite $G$: correlation decay at weak coupling~\cite{Adhikari2025}; 3D $U(1)$, Villain action: area law at all couplings~\cite{Gopfert1982}.) \\
\bottomrule
\end{tabularx}
\caption{Status of the problem and its relatives (6 November 2026, with the 10 November release noted).}
\label{tab:status}
\end{table}

%=====================================================================
\section{A candidate formal statement in Lean 4}\label{sec:statement}
%=====================================================================

The statement is a Lean~4~\cite{Moura2021} project built on Mathlib~\cite{Mathlib2020}
(toolchain and Mathlib tag \lean{v4.36.0-rc2}, Mathlib commit \lean{9d581df}). It imports only
Mathlib, contains definitions and proofs only, and confines every \lean{sorry} to one file of
open problems. About 1{,}700 lines are spread over \lean{OS.lean} (axioms),
\lean{Lattice.lean} (lattice gauge theory), \lean{SU.lean} ($SU(N)$),
\lean{Statement.lean} (targets), \lean{OpenProblems.lean}, \lean{Sanity.lean},
\lean{NonVacuity/*.lean} and \lean{Audit.lean}.

\subsection{Design choices}
\begin{enumerate}[leftmargin=*]
\item \textbf{Euclidean OS form.} OS form is closest to every constructive route; OS~II yields a
  Wightman theory.
\item \textbf{Multilinear Schwinger functions.} For fields $\varphi_i$, $i\in\iota$, transforming
  under a real representation $\rho$ of $O(4)$, the $n$-point function is a continuous
  multilinear form
  $S_n(f_1,\dots,f_n)=\langle\varphi(f_1)\cdots\varphi(f_n)\rangle$ on
  $f_k\in\Sch(\R^4,\C^\iota)$. This avoids building tensor products of Schwartz functions.
\item \textbf{Disjoint supports.} Composite-operator correlations have scheme-dependent contact
  terms. Following OS's use of $\Sch_0$, every axiom except growth is imposed only on tuples
  with pairwise disjoint supports; reflection positivity is imposed on products supported in
  ordered positive-time slabs, the multilinear analogue of OS's $\Sch_+$.
\item \textbf{Constructed operations, not relations.} Euclidean motions, translations and the
  reflection-conjugation $\Theta$ are built as continuous linear maps from Mathlib's
  \lean{SchwartzMap\dotbr compCLMOfContinuousLinearEquiv}, \lean{compSubConstCLM} and
  \lean{postcompCLM}. A purely relational definition would risk vacuity if no related function
  existed.
\end{enumerate}

\subsection{Axioms}
Write $\Theta f=\overline{\rho(\theta)^{\mathsf T}f(\theta\,\cdot)}$ with $\theta$ time reflection, and for
finite linear combinations $F=\sum_j c_j\,\varphi(f_{j,1})\cdots\varphi(f_{j,n_j})$ and
$G=\sum_l d_l\,\varphi(g_{l,1})\cdots\varphi(g_{l,m_l})$ of ordered-positive monomials define the
reflection-positivity form
\[
B(F,G)=\sum_{j,l}\overline{c_j}\,d_l\;S_{n_j+m_l}\big(\Theta f_{j,n_j},\dots,\Theta f_{j,1},
g_{l,1},\dots,g_{l,m_l}\big).
\]
\lean{OSAxioms} consists of:
\begin{itemize}[leftmargin=*]
\item normalization ($S_0=1$);
\item an equivalent form of \eqref{eq:E0prime}: for each $n$, $S_n$ agrees on disjointly supported
  products with a distribution $T_n$ on $\Sch((\R^4)^n)$ obeying
  $|T_n(F)|\le\alpha(n!)^\beta|F|_{ns}$ (equivalent to OS~II's condition on $\Sch_0$ by
  Hahn--Banach and density of disjointly supported products, with norm-convention constants
  absorbed into $(n!)^\beta$; this equivalence is argued, not formalized);
\item reality;
\item Euclidean covariance under $f\mapsto\rho(R^{-1})^{\mathsf T}f(R^{-1}(\cdot-a))$ for all
  $R\in O(4)$, $a\in\R^4$;
\item permutation symmetry;
\item reflection positivity $B(F,F)\ge0$;
\item the cluster property
  $S(\Theta A^*\cdot T_{\lambda a}B)\to S(\Theta A^*)S(B)$ as $\lambda\to\infty$, for every
  \emph{non-zero} spatial $a$.
\end{itemize}

\subsection{The mass gap}
With the truncated form $B_T(F,G)=B(F,G)-B(F,1)B(1,G)$ and Euclidean time translation $T_t$:

\begin{definition}[\lean{HasMassGap}]
$\Delta>0$ and, for all ordered-positive $F,G$ and all $t\ge0$,
$|B_T(F,T_tG)|\le e^{-\Delta t}\,B_T(F,F)^{1/2}B_T(G,G)^{1/2}$.
\end{definition}

\begin{proposition}
After OS reconstruction, \lean{HasMassGap}~$\Delta$ is equivalent to
$\|e^{-tH}(1-P_\Omega)\|\le e^{-\Delta t}$ for $t\ge0$, i.e.\ the spectrum of the Hamiltonian
on $\Omega^\perp$ lies in $[\Delta,\infty)$ and the vacuum is unique.
\end{proposition}
\begin{proof}[Sketch]
Reconstruction maps ordered-positive $F$ to vectors $\hat F$ in the OS Hilbert space $\mathcal H$
generated by the fields, with $\langle\hat F,\hat G\rangle=B(F,G)$ and $\widehat{T_tG}=e^{-tH}\hat G$,
$H=H^*\ge0$. These vectors are dense because finite sums of products of slab-supported test
functions are dense in OS's $\Sch_+$; this density step, like the proposition itself, is argued but
not yet formalized. Then $B_T$ is the inner
product of the projections onto $\Omega^\perp$, and the inequality on a dense set is the stated
operator bound. The spectral theorem converts the operator bound into the spectral statement.
\end{proof}
\noindent With Lorentz covariance (from OS~II), this is the Jaffe--Witten gap. The finiteness of
$m$ follows from non-triviality.

\subsection{Identifying the theory as Yang--Mills}
The OS axioms, non-triviality and a gap are satisfied by theories that are not Yang--Mills (for
example Wick powers of a massive free field). Four clauses pin the theory down.

\paragraph{Non-triviality (\lean{NonTrivial}).}
The fields are centred; some two-point function and some fourth cumulant are non-zero at
separated points; and some state orthogonal to the vacuum has positive norm.

\paragraph{The lattice clause (\lean{IsWilsonLimit}).}
For a compact group $G$ with a faithful continuous unitary representation $\pi:G\to U(d)$,
consider the periodic lattice $(\mathbb{Z}/(L+1))^4$, link variables $U$, plaquettes
$U_{\mu\nu}(x)$, and Wilson's action~\cite{Wilson1974}
\[
S_\beta(U)=\beta\sum_x\sum_{\mu<\nu}\Big(1-\tfrac1d\,\mathrm{Re}\,\tr\,\pi(U_{\mu\nu}(x))\Big),
\]
with the Gibbs measure built from normalized Haar measure. The lattice field strength
$F_{\mu\nu}(x)$ is the traceless Hermitian part of $(\pi(U_{\mu\nu})-\pi(U_{\mu\nu})^\dagger)/2i$,
and $T_{\mu\nu\kappa\lambda}(x)=\mathrm{Re}\,\tr F_{\mu\nu}(x)F_{\kappa\lambda}(x)$. Smeared
fields are Riemann sums $a^4\sum_x f(ax)\,Z\,(T(x)-\langle T(x)\rangle)$ with renormalization
matrices $Z$ valued in the 21-dimensional space of tensors with the symmetries of
$\tr F_{\mu\nu}F_{\kappa\lambda}$. $S$ is a \emph{Wilson limit} if, along some sequence with
$a_k\to0$, $a_k(L_k+1)\to\infty$ and $\beta_k\to\infty$, the lattice Schwinger functions
converge to $S$ on compactly supported test functions with disjoint supports. O(4) acts on the
fields by $R\otimes R\otimes R\otimes R$; the Lean representation is a monoid homomorphism built
from Kronecker products.

\paragraph{Full field content (\lean{FullFieldContent}).}
Every non-zero direction $v$ in the 21-dimensional space yields a non-zero state
$\varphi_v(h)\Omega$. Without this clause, the existential renormalization matrices could project
the theory onto a single channel such as $\tr F^2$; this is a finding of the review in
Section~\ref{sec:review}.

\paragraph{Short-distance asymptotic freedom (\lean{ShortDistanceAF}).}
In the scalar channel ($v_{\mu\nu\kappa\lambda}=\delta_{\mu\kappa}\delta_{\nu\lambda}-\delta_{\mu\lambda}\delta_{\nu\kappa}$),
for compactly supported $h_1,h_2$ with disjoint supports and dilations
$D_\lambda h(x)=\lambda^{-4}h(x/\lambda)$, we require for some $C>0$
\[
\lambda^8\log^2(1/\lambda)\;S_2\big(D_\lambda h_1,D_\lambda h_2\big)\;\longrightarrow\;
C\iint h_1(x)\,h_2(y)\,|x-y|^{-8}\,dx\,dy\qquad(\lambda\to0^+).
\]
The motivation is perturbative. The renormalization-group-invariant trace-anomaly operator
$\Theta=(\beta(g)/2g)F^a_{\mu\nu}F^a_{\mu\nu}$, with $\beta(g)=\mu\,dg/d\mu=-b_0g^3-b_1g^5+\cdots$~\cite{Collins1977},
has at leading order correlator $\propto\bar g^4(1/|x|)\,|x|^{-8}$, where $\bar g(\mu)$ is the
running coupling at scale $\mu$; with $\bar g^2(\mu)\sim 1/(2b_0\log(\mu/\Lambda))$ this gives a
$\log^{-2}$ factor absent in free field theory. The clause therefore excludes free massless gluons even without a gap. It is a
\emph{requirement}; its validity for Yang--Mills is expected from perturbation theory but is
not proved.

\paragraph{Asymptotic scaling (\lean{AsymptoticScaling}, $SU(N)$ only).}
With $\beta=2N/g_0^2$, $b_0=11N/(3(4\pi)^2)$~\cite{Gross1973,Politzer1973} and
$b_1=34N^2/(3(4\pi)^4)$~\cite{Caswell1974,Jones1974},
$\beta_k/(2N)-2b_0\log(1/a_k)-(b_1/b_0)\log\log(1/a_k)$ converges as $k\to\infty$.

\subsection{The targets}
\begin{Verbatim}
def ClayYangMills : Prop :=
  ∀ (d₀ : ℕ) (G : Type) [Group G] [TopologicalSpace G] [IsTopologicalGroup G]
    [CompactSpace G] [ConnectedSpace G] [MeasurableSpace G] [BorelSpace G]
    (P₀ : GaugeGroup d₀ G),
    IsSimpleMatrixLieAlgebra P₀.lieAlgebra → ExistsWithMassGapFor G

def IsYangMillsTheoryAlong P S a β : Prop :=
  OSAxioms tensorRep S ∧ NonTrivial tensorRep S ∧ FullFieldContent S ∧
    ShortDistanceAF S ∧ IsWilsonLimit P S a β
\end{Verbatim}
\lean{ExistsWithMassGapFor G} asserts that, for some faithful unitary representation of $G$,
there is a Yang--Mills theory in this sense with \lean{HasMassGap}~$\Delta$ for some $\Delta$.
\lean{ClayExistence} drops the gap. \lean{SUNExists N} uses the fundamental representation of
$SU(N)$ and adds \lean{AsymptoticScaling}. \lean{SUNExistsWithMassGap N} uses the fundamental
Wilson action and requires the gap. The docstrings of \lean{SUNExists} and
\lean{su\_yang\_mills\_exists} call it the shape or scope of OpenAI's announced result. That wording
predates the 10 November release, was not checked against OpenAI's Lean code, and is left unchanged
because the statement-file hashes are fixed (Section~\ref{sec:data}).

\subsection{Milestone ladder}
\lean{SUNExists} bundles two genuinely infrared clauses, full $O(4)$ covariance and the cluster
property. Given the other axioms, the cluster property is equivalent to uniqueness of the vacuum;
since \lean{IsWilsonLimit} uses periodic tori with no boundary conditions, the periodic limit itself
must have a unique vacuum. We therefore define weaker rungs:
\lean{SUNExistsM7a} (hypercubic covariance, no cluster clause) and \lean{SUNExistsM7b} (full
$O(4)$, no cluster clause), and prove \lean{SUNExists}~$\to$~M7b~$\to$~M7a in Lean. Only
\lean{SUNExists} (and \lean{ClayExistence} for general $G$) corresponds to the existence part of
the Clay problem.

\subsection{Deviations from the Clay text}
The Clay text is informal in places, so no formal statement can be proved equivalent to it.
Table~\ref{tab:dev} records each deviation and its direction.

\begin{table}[ht]
\centering\small
\begin{tabularx}{\textwidth}{@{}p{4.3cm}p{2.7cm}X@{}}
\toprule
Aspect & Direction & Consequence \\
\midrule
Gap only on the cyclic space of $\tr F F$ (a subspace of the C-even sector) & weaker & a full proof must add C-odd and higher-dimension fields \\
Wilson-lattice origin required & stronger/different & other regulators need a universality theorem; Douglas~\cite{Douglas2004} notes the lattice is ``not essential'' \\
AF fingerprint in the scalar two-point function only & weaker & no stress-tensor or OPE clause \\
Subsequential limits accepted & choice & axioms and gap are required of the limit, so JW's ``no weak existence'' is met \\
E0$'$ rather than Glimm--Jaffe OS1 & OS~II's own condition & OS1 is incompatible with the expected singularity of $\tr F^2$ (our argument; see Section~\ref{sec:problem}) \\
Representation chosen existentially & equivalent intent & a fundamental-representation variant is kept separately \\
\bottomrule
\end{tabularx}
\caption{Deviations of the candidate statement from the Clay problem description.}
\label{tab:dev}
\end{table}

%=====================================================================
\section{Machine-checked results}\label{sec:checked}
%=====================================================================

Table~\ref{tab:checked} lists the proved theorems. For each, \lean{\#print axioms} reports only
\lean{propext}, \lean{Classical.choice} and \lean{Quot.sound}. The open problems additionally
depend on \lean{sorryAx}.

\begin{table}[ht]
\centering\small
\begin{tabularx}{\textwidth}{@{}lX@{}}
\toprule
Theorem & Content \\
\midrule
\lean{trivial\_osAxioms}, \lean{trivial\_massGap} & The trivial theory ($S_0=1$, $S_n=0$) satisfies every OS axiom and has a gap: the axioms as formalized are consistent. \\
\lean{trivial\_not\_nonTrivial} & \ldots\ but it is not \lean{NonTrivial}: the clause does real work. \\
\lean{exists\_spatial\_ne\_zero} & Regression guard: the cluster axiom quantifies over a non-empty set. \\
\lean{partitionFn\_pos}, \lean{expect\_const} & For second-countable $G$ (which includes every compact group with a faithful continuous unitary representation), the Wilson partition function is positive, so lattice expectations are genuine averages; $\langle c\rangle=c$. \\
\lean{isCompact\_specialUnitaryGroup} & $SU(N)$ is compact (not previously in Mathlib). \\
\lean{su\_connected} & $SU(N)$ is connected (via \lean{su\_pathConnectedSpace}): $U(N)$ is shown path-connected from Mathlib's \lean{Unitary.joined}, and a determinant-correcting retraction maps it onto $SU(N)$. \\
\lean{suGauge\_lieAlgebra} & The Lie algebra $\{X:\ e^{tX}\in SU(N)\ \forall t\}$ equals $\mathfrak{su}(N)$; uses $\det e^{Y}=e^{\tr Y}$ for skew-Hermitian $Y$, proved via the spectral theorem. \\
\lean{su\_isSimple} & $\mathfrak{su}(N)$, $N\ge2$, is a simple real Lie algebra. \\
\lean{su\_mass\_gap\_of\_clay} & \lean{ClayYangMills} implies the conclusion for $SU(N)$: the general statement is not vacuous. \\
\lean{SUNExists.toM7b}, \lean{SUNExistsM7b.toM7a} & The milestone ladder is ordered. \\
\bottomrule
\end{tabularx}
\caption{Fully proved results in the statement package.}
\label{tab:checked}
\end{table}

\noindent Open (\lean{sorry}): \lean{clay\_yang\_mills}, \lean{clay\_existence},
\lean{su\_yang\_mills\_exists}, \lean{su\_yang\_mills\_\allowbreak mass\_gap}. One classical but unformalized
obligation, \lean{os\_nonTrivial\_\allowbreak fullField\_consistent}, asserts that the OS axioms are
consistent with \lean{NonTrivial} and \lean{FullFieldContent} (witness: Wick-ordered composites
of the free Maxwell field with values in $\mathfrak{su}(N)$). Consistency of the full Yang--Mills
conjunction, including \lean{ShortDistanceAF} and the lattice clause, is the open problem itself.

%=====================================================================
\section{Adversarial review of the statement}\label{sec:review}
%=====================================================================

\subsection{Method}
Version 0.1 of the statement was reviewed by six AI reviewer agents (separately prompted,
context-isolated instances of the same model), each with a distinct lens:
\begin{itemize}[leftmargin=*]
\item fidelity to OS~I/II;
\item mass gap and non-triviality;
\item whether the lattice clause pins down Yang--Mills;
\item Lean technicalities and junk values;
\item comparison with the Clay text and with public formalizations~\cite{Cherkis2026,Yin2026,OpenAI2026m,OpenAI2026j};
\item a red team instructed to prove or disprove the open theorems cheaply.
\end{itemize}
Reviewers worked read-only on the statement but could compile scratch Lean files, so exploits
were demonstrated, not argued. Each finding then went to a separate verifier agent (same model)
instructed to refute it, defaulting to ``refuted'' when in doubt. Of 14 verified findings, 13
survived and 1 was refuted. Compiling exploits in Lean, together with the non-vacuity witnesses of
Section~\ref{sec:checked}, is in the spirit of the mechanical vacuity checks of Mei and
Douglas~\cite{Mei2026}, who combine LLM-based semantic review with mechanical numerical and vacuity
checks and report that LLM judges miss errors those checks catch; our verification of findings, by
contrast, was LLM-judged.

\subsection{Findings}
\begin{itemize}[leftmargin=*]
\item \textbf{Critical: the cluster axiom admitted $a=0$.}
  \begin{itemize}
  \item Found separately by three reviewer agents (same model, different lenses).
  \item The condition ``$a$ spatial'' was encoded as $a^0=0$, which the zero vector satisfies.
    At $a=0$ the sequence is constant, so the axiom forced exact factorization
    $S(\Theta A^*B)=S(\Theta A^*)S(B)$. That collapses the OS Hilbert space to $\C\Omega$.
  \item Consequences: \lean{HasMassGap} followed from the axioms, and the axioms together with
    \lean{NonTrivial} were unsatisfiable. The open theorems were therefore \emph{provably false},
    and reviewers' Lean files proved the collapse.
  \item Fixed by requiring $a\neq0$. The reviewers' exploit files (kept in \lean{scratch/},
    outside the build) no longer compile, and \lean{exists\_spatial\_ne\_zero} shows that the
    corrected axiom is not vacuous; no automated regression test yet guards the fixes.
  \item The trivial-theory sanity check could not have caught this, because the trivial theory
    factorizes exactly. Non-degenerate witnesses are required.
  \end{itemize}
\item \textbf{Major: growth weaker than E0$'$.} The first version bounded $S_n$ by products of
  per-factor Schwartz norms of order $n\cdot s$. That bound follows from E0$'$ but is not known to
  imply it: the product-to-E0$'$ argument in the appendix to OS~II needs per-factor norms of fixed
  order. So OS~II's reconstruction theorem could not be invoked. Replaced by E0$'$ on configuration
  space.
\item \textbf{Major: single-channel collapse.} Unconstrained renormalization matrices let a
  witness keep only $\tr F^2$. Lean lemmas showed that the Wilson-limit, reflection-positivity
  and gap clauses all survive such a projection. Fixed by \lean{FullFieldContent} and a range
  condition on $Z$.
\item \textbf{Major: missing short-distance requirement.} The JW requirement of agreement with
  asymptotic freedom had no formal counterpart. Added \lean{ShortDistanceAF}, a scalar two-point
  fingerprint. The stress-tensor and OPE parts of the requirement remain unformalized and are
  recorded in Table~\ref{tab:dev}; this finding is therefore also only partly fixed.
\item \textbf{Major: method and representation dependence.} Requiring a Wilson-lattice limit makes
  the statement method-dependent and, in that respect, stronger than the Clay text, and universal
  quantification over faithful representations built lattice universality into the target. The
  representation was made existential. The Wilson-lattice requirement remains and is recorded in
  Table~\ref{tab:dev}; a method-independent variant is future work. This finding is therefore only
  partly fixed.
\item \textbf{Minor: documentation.} The renormalization matrices were described as ``finite''.
  Any non-trivial limit is expected to need $\|Z_k\|\to\infty$, like $a_k^{-4}$ times
  channel-dependent powers of $\log(1/a_k)$ (perturbatively $a_k^{-4}$ for the trace-anomaly scalar
  and $a_k^{-4}g_0^{-2}$ for the stress tensor).
\end{itemize}

%=====================================================================
\section{Research landscape}\label{sec:landscape}
%=====================================================================

We summarise the eleven library notes. Each claim there carries a label: established theorem,
claimed but unverified, heuristic, disputed or open. The notes also label some claimed solutions as
containing suspected errors; these are AI agents' assessments, not reviewed by any human expert,
and should not be read as refereed judgements.

\paragraph{Constructive QFT and renormalization.}
\begin{itemize}[leftmargin=*]
\item Superrenormalizable models were constructed in the 1970s~\cite{Glimm1987,Feldman1976}.
\item The massive Gross--Neveu model in two dimensions is generally cited as the first non-planar,
  non-hierarchical asymptotically free renormalizable theory to be
  constructed~\cite{Gawedzki1985a,Feldman1986}; planar $\varphi^4_4$ with the wrong sign was
  constructed earlier~\cite{Hooft1983,Rivasseau1984}.
\item Modern renormalization-group methods include~\cite{Bauerschmidt2019,Brydges2015a}, and
  cluster expansions~\cite{Kotecky1986}.
\item Aizenman and Duminil-Copin proved that scaling limits of critical four-dimensional Ising and
  (single-component, nearest-neighbour) lattice $\varphi^4_4$ models are Gaussian~\cite{Aizenman2021}.
  With the one-loop analysis of Coleman and Gross~\cite{Coleman1973}, this supports, but does not
  prove, the expectation that a non-trivial four-dimensional theory must be asymptotically free.
\end{itemize}

\paragraph{Ba{\l}aban's programme.}
In 1984--1989 Ba{\l}aban proved ultraviolet stability bounds (an upper bound on the partition
function, and bounds on the effective densities) for lattice Yang--Mills on 3D tori, and on 4D tori
conditionally on the running couplings staying small, using a gauge-covariant block-spin
renormalization group. The work spans about 400 pages of
\emph{Commun.\ Math.\ Phys.}~\cite{Balaban1984a,Balaban1984b,Balaban1985a,Balaban1985c,Balaban1985e,Balaban1985f,Balaban1985g,Balaban1987,Balaban1988b,Balaban1988d,Balaban1989b,Balaban1989c}.
\begin{itemize}[leftmargin=*]
\item The 4D theorem is conditional on the running coupling staying small, and an announced
  companion result on the coupling flow was, according to secondary sources~\cite{Dimock2013a}, never
  published.
  Checking this against the primary text is our first stop rule (SR1).
\item According to these secondary accounts, correlation functions, convergence, infinite volume,
  the OS axioms and the gap are not treated in Ba{\l}aban's work.
\item Dimock's expositions~\cite{Dimock2013a,Dimock2013b,Dimock2013c,Dimock2014,Dimock2022} are the
  only sustained, refereed modern re-derivation of Ba{\l}aban's method; they treat $\varphi^4_3$ and
  QED$_3$, not non-Abelian Yang--Mills. The 2026 AI-generated manuscripts~\cite{Sussman2026,Aronow2026}
  are unrefereed.
\end{itemize}

\paragraph{Rigorous lattice results.}
\begin{itemize}[leftmargin=*]
\item Reflection positivity of Wilson's action, and confinement at strong coupling, are due to
  Osterwalder and Seiler~\cite{Osterwalder1978,Seiler1982}.
\item Large-$N$ and strong-coupling analyses~\cite{Chatterjee2019,Basu2018,Jafarov2016}, and exact
  random-surface representations of Wilson-loop expectations~\cite{Cao2025c}.
\item Finite gauge groups~\cite{Cao2020,Chatterjee2020,Forsstrom2022b}, including correlation
  decay at weak coupling~\cite{Adhikari2025}.
\item Stochastic-analysis gaps at strong coupling~\cite{Shen2023}, and area-law
  regimes~\cite{Cao2025a,Cao2025b}.
\item Conditional reductions of confinement to correlation decay~\cite{Durhuus1980,Chatterjee2021}.
\item None reaches the weak-coupling regime of the continuum limit in four dimensions.
\end{itemize}

\paragraph{Stochastic quantization.}
\begin{itemize}[leftmargin=*]
\item Langevin dynamics, gauge-covariant state spaces and universality are established in two
  dimensions~\cite{Chandra2022,Chevyrev2019,Chevyrev2026b}.
\item In three dimensions, local solution theory and uniqueness of renormalization are
  established~\cite{Chandra2024b,Chevyrev2026d}, and Cao and Chatterjee~\cite{Cao2023,Cao2024a}
  supply a gradient-flow state space.
\item A construction of the non-Abelian Yang--Mills--Higgs measure on the two-dimensional torus,
  for all fixed couplings, was posted on 2 November 2026~\cite{Bringmann2026b}. It is an
  unrefereed preprint, and its authors disclose that several lemmas were proved by AI models.
\item In four dimensions the Yang--Mills Langevin equation is scaling-critical, which places it
  outside the locally subcritical class to which regularity structures~\cite{Hairer2014} and
  related methods apply.
\item Gradient-flow observables~\cite{Luscher2010a,Luscher2011}, which are UV-finite to all
  orders of perturbation theory, are a natural observable class.
\end{itemize}

\paragraph{Mass-gap mechanisms and sanity filters.}
Any purported proof of a 4D gap must:
\begin{itemize}[leftmargin=*]
\item fail for 4D $U(1)$, which has a massless Coulomb phase~\cite{Guth1980,Frohlich1982a};
\item not prove confinement for finite groups;
\item respect the $\theta=\pi$ anomaly~\cite{Gaiotto2017}.
\end{itemize}
Other relevant results:
\begin{itemize}[leftmargin=*]
\item Rigorous confinement (an area law at all couplings) holds for 3D $U(1)$ lattice gauge theory
  with the Villain action~\cite{Gopfert1982}.
\item Controlled (physics-level, non-rigorous) semiclassical analyses exhibit gapped, confining
  regimes on compactified geometries~\cite{Tanizaki2022}; their connection to $\R^4$ rests on an
  adiabatic-continuity conjecture.
\item Finite-size uniqueness criteria exist~\cite{Dobrushin1985}.
\item Computer-assisted renormalization group~\cite{Ebel2025} is a sober calibration of what
  certified numerics can currently reach.
\item Lattice Monte Carlo studies~\cite{Morningstar1999,Lucini2013} give strong numerical
  evidence for the gap.
\end{itemize}

\paragraph{AI-era claims (all unrefereed).}
\begin{itemize}[leftmargin=*]
\item OpenAI's announced results and repositories~\cite{OpenAI2026j,OpenAI2026m,OpenAI2026n}.
  These include a claimed exponential decay of spin correlations at every positive temperature for
  the 2D nearest-neighbour $O(n)$ model with $n\ge3$, with a Lean development covering finite
  free-boundary subgraphs~\cite{OpenAI2026b}.
\item Sussman--Alp\"oge's 3D manuscripts~\cite{Sussman2026}.
\item Aronow--Lopatto's 3D continuum limit~\cite{Aronow2026}.
\item Craig's gauge-invariant lattice measure for chiral fermions~\cite{Craig2026a}.
\item Public calls for release of OpenAI's material~\cite{AGMAI2026}.
\end{itemize}

%=====================================================================
\section{A research programme}\label{sec:programme}
%=====================================================================

\subsection{What the premise implies}
If OpenAI's claim is correct and, as reported, its proof is built on Wilson's lattice and
Ba{\l}aban's renormalization group and covers fermions, then:
\begin{itemize}[leftmargin=*]
\item \emph{its} existence theorem is true;
\item a Ba{\l}aban-type renormalization group on Wilson's lattice can be carried to infinite
  volume without a mass gap;
\item the pipeline handles fermions.
\end{itemize}
The premise does \emph{not} imply that our \lean{SUNExists} is provable without new infrared
mathematics. The cluster clause requires vacuum uniqueness for periodic limits. Limits of periodic
tori can be non-trivial mixtures (as for the low-temperature Ising model, where every periodic limit
is $\tfrac12(\mu^++\mu^-)$), and the ergodic components of such a mixture need not be periodic
limits; ruling this out for Yang--Mills needs infrared input. Nor does the premise say anything
about the gap.

\subsection{Targets and tracks}
The programme climbs a ladder of frozen Lean statements, each adversarially reviewed and
hash-recorded before proof search:
\begin{itemize}[leftmargin=*]
\item finite-lattice foundations;
\item a Gaussian $U(1)_4$ calibration rung;
\item a 4D bounded-field multiscale renormalization group carrying the running coupling and a
  $\Theta$-source;
\item a conditional milestone in which Ba{\l}aban's published theorems enter as explicitly named
  hypotheses, as in OpenAI's repositories;
\item finite-volume continuum limits;
\item $T^3_\ell\times\R$ at small $\ell$, the first genuine Hilbert space and gap;
\item the rungs M7a/M7b/M7c.
\end{itemize}
The work is organized in five tracks:
\begin{itemize}[leftmargin=*]
\item \textbf{A:} statement, bridges and release audit;
\item \textbf{B:} demand-driven foundations in Lean;
\item \textbf{C:} the ultraviolet analysis following Ba{\l}aban, starting from primary sources;
\item \textbf{D:} infinite volume without a gap;
\item \textbf{E:} the mass gap, as a capped research portfolio.
\end{itemize}
Within Track~E, a desk analysis identified an obstruction. A finite-size criterion that is
uniform over all volumes and boundary conditions cannot hold in four dimensions beyond bare
coupling of order one. Take the time-like $SU(2)$ links of a height-one slab and freeze all other
links to the identity; the conditional measure is then exactly the 3D $O(4)$ model at the same
$\beta$, which has long-range order for large $\beta$. This is our heuristic argument, not a
theorem (with generic frozen links one gets an $O(4)$ model in an $SO(4)$ background, which can be
frustrated). It directs any certificate approach to cube-restricted criteria.

\subsection{AI operating model}
The design draws on the published record of OpenAI's Navier--Stokes swarm. That swarm reportedly
used groups of agents with in-group messaging, a seeding stage from an Euler result, consolidation
of insights across groups, and a separate formalization
stage~\cite{OpenAI2026g,Howlett2026h,OpenAI2026k}. OpenAI researchers credit base-model strength
over multi-agent organisation~\cite{Patel2026}, and critics have documented the risks of scope
loopholes~\cite{Howlett2026b,Constantin2026c} and of Lean/paper drift~\cite{Bastounis2026}. The
design also draws on OpenAI's Lean repositories: Comparator-gated challenges, honest-axiom
patterns, conditional hypotheses, and metadata
files~\cite{OpenAI2026m,OpenAI2026j,OpenAI2026l,LeanFRO2026,Mathlib2026b}.
Its principles are:
\begin{enumerate}[leftmargin=*]
\item \textbf{The compiled Lean skeleton is the coordination substrate.} Each rung theorem is
  proved in Lean from a small number of named module contracts.
\item \textbf{Agents read contracts, not proofs.} Context is bounded.
\item \textbf{Mechanical checks first, model judgement second, the kernel last.} LLM skeptics
  are a rejection filter calibrated by planted errors, never an acceptance gate.
\item \textbf{Swarms only where the decomposition is unknown.}
\item \textbf{Provenance by construction.}
\end{enumerate}
Blueprint tooling~\cite{Massot2026,Becker2024} and Comparator~\cite{LeanFRO2026} provide the
gates.

\subsection{Honest odds and stop rules}
The programme adopts the subjective estimates of its feasibility review. At a sustained budget
of about $10^{10}$ output tokens per year:
\begin{itemize}[leftmargin=*]
\item an independent, Comparator-checked \lean{SUNExists} is $\le1\%$ within 12 months and
  3--8\% within 36 months;
\item a Lean bridge from OpenAI's theorem to ours is 60--80\% likely within 3 months
  \emph{if} their statement implies ours up to bookkeeping;
\item the 4D mass gap is $<0.1\%$ within 12 months and $<0.5\%$ within 36 months at this budget,
  and $<1\%$ within 36 months even at a Navier--Stokes-scale budget of about $10^{11}$ tokens per
  year.
\end{itemize}
Nine stop rules tie continuation to measurable checkpoints:
\begin{itemize}[leftmargin=*]
\item choosing the clean-room/bridge/consume mix fourteen days after OpenAI's release;
\item confirmation of the Ba{\l}aban gap in the primary text;
\item calibrated derivation of $b_0$ in Ba{\l}aban's scheme;
\item measured closure rate on a pilot formalization of Dimock's $\varphi^4_3$ analysis;
\item canary recall of the review gates;
\item spend caps;
\item Track~E caps;
\item switching to verification if a complete public 4D manuscript appears;
\item withdrawing the clean-room label from artefacts affected by a firewall breach.
\end{itemize}

\subsection{Independence and ethics}
Until OpenAI's release (10 November 2026), all work was clean-room by necessity; by policy it stays
clean-room for fourteen days after the release (to 24 November), after which the author chooses among three modes:
\begin{itemize}[leftmargin=*]
\item \emph{clean-room}: no artefact derived from the release;
\item \emph{bridge}: our statement proved from theirs;
\item \emph{consume}: study the release and re-derive it.
\end{itemize}
Leakage controls and labelling rules accompany each mode. The programme commits to:
\begin{itemize}[leftmargin=*]
\item contacting Sussman and Alp\"oge before duplicating their four-dimensional modules;
\item respecting the committee's process;
\item disclosing AI use;
\item publishing any claimed proof, or any formal result beyond this candidate statement, only
  after Comparator checking and review by human experts, with no press releases. This paper itself
  is released without human expert review (Section~\ref{sec:ai}), as a candidate statement and a
  plan, not as a result.
\end{itemize}

%=====================================================================
\section{Limitations}
%=====================================================================
\begin{itemize}[leftmargin=*]
\item The statement is a \emph{candidate}. It has been reviewed adversarially by AI agents, not
  audited by human experts in constructive QFT. Its fidelity to the Clay problem is argued, not
  proved; Table~\ref{tab:dev} lists the known deviations.
\item The mass-gap clause covers only the cyclic space of the dimension-four fields.
\item Consistency of OS + \lean{NonTrivial} + \lean{FullFieldContent} is asserted but not yet
  formalized.
\item The library relies on secondary sources in important places. No paper of Ba{\l}aban's series
  was read in full (the publisher's site blocked automated access), so
  Section~\ref{sec:landscape}'s account of his work rests on Dimock's expositions, other secondary
  accounts and an unvetted third-party transcription. OpenAI's blog posts~\cite{OpenAI2026g,OpenAI2026p}
  returned HTTP~403 and were reconstructed from press and blog reports. Douglas's 2026
  review~\cite{Douglas2026f} was not read. Bibliography entries that could not be confirmed are
  flagged.
\item The Proposition of Section~3.3 and the equivalence of the multilinear axioms with OS's
  formulation are argued, not formalized.
\item The research programme's odds are subjective. Its operating model is untested at the
  scale required.
\item Nothing here independently evaluates OpenAI's claimed proof. Its Lean code, released on
  10 November 2026~\cite{Douglas2026g}, was not examined for this version; a comparison of the two
  statements is left to a later version.
\end{itemize}

%=====================================================================
\section{Disclosure of AI use}\label{sec:ai}
%=====================================================================

This disclosure follows, as far as they apply to an individual's project, the disclosure items
that the Advisory Group on Mathematics and AI lists for AI-generated mathematics~\cite{AGMAI2026a}:
model, prompts, time, compute, and how AI was used. Those guidelines are addressed mainly to AI
laboratories. Provided here: the model family and environment; the prompt given to every subagent
(in the workflow scripts, \lean{dev/workflows/}); and wall-clock times, tool calls and subagent
token counts. Not included in this deposit: exact model versions (the Git commit messages in the GitHub repository carry an automatic
co-author line naming only the model of the committing session), the author's natural-language
instructions to the orchestrating session, that session's transcript and token count, summaries of
the models' reasoning, and monetary cost.

\paragraph{Systems used.}
The work was carried out with Anthropic's Claude large language models, operated through the Claude
Code agent environment. That environment includes a workflow facility that orchestrates many
subagents in parallel, file, shell, web-search and web-fetch tools, and the Lean toolchain. Research
subagents also used the alphaXiv paper-search and PDF question-answering service, a third-party tool
whose answers may themselves be machine-generated. The models' training data predate the September
2026 events discussed here; all information about those events was obtained by live retrieval
during the work.

\paragraph{What the human author did.}
The author:
\begin{itemize}[leftmargin=*]
\item set the goal and the working premise;
\item directed the work in natural language and answered the agents' questions;
\item chose the repository visibility, publication route and licences;
\item takes responsibility for publication.
\end{itemize}
The author did not hand-write the mathematics, the Lean code or the prose.

\paragraph{Human review.}
The author read the paper before deposit, including the passages revised on 10--11 November, but is
not trained in constructive quantum field theory or
in Lean, and could not independently check the mathematics, the physics or the Lean code. For these
the paper relies on the Lean kernel (Section~\ref{sec:checked}) and on the AI review described in
Section~\ref{sec:review}. No human domain expert has reviewed this work. No text-overlap
(plagiarism) check was run on the AI-generated text; quotations were checked against their sources
by AI agents only. Corrections are welcome as issues on the GitHub repository.

\paragraph{What AI systems did.}
AI agents performed essentially all of the following:
\begin{itemize}[leftmargin=*]
\item literature search and summarization;
\item drafting of the library notes, the problem description, the plan and this paper (the
  orchestrating session wrote the Lean statement, its version~0.2 fixes, the problem description,
  the plan and this paper; subagents wrote the notes, reviews, critiques and non-vacuity proofs);
\item design and writing of the Lean statement and all Lean proofs;
\item adversarial review and verification of the statement, the notes, the plan and this paper.
\end{itemize}

\paragraph{Scale and structure (subagents only).}
\begin{itemize}[leftmargin=*]
\item \textbf{Research library.} 24 subagents (11 researchers, 11 adversarial citation and claim
  verifiers, an indexer and a completeness critic): about $9.1\times10^6$ tokens, 2{,}560 tool
  calls, about 2.0~hours.
\item \textbf{Statement review.} 20 subagents (6 reviewers, 14 verifiers): about $3.6\times10^6$
  tokens, 272 tool calls, about 24~minutes.
\item \textbf{Plan critique.} 5 subagents: about $2.0\times10^6$ tokens, 194 tool calls, about
  28~minutes.
\item \textbf{Non-vacuity Lean proofs.} 3 subagents: about $0.6\times10^6$ tokens, 70 tool calls,
  about 12~minutes.
\item \textbf{Paper verification.} 7 subagents (four citation checkers, a factual-claims auditor, a
  mathematics reviewer and a disclosure reviewer): about $2.0\times10^6$ tokens, 336 tool calls,
  about 22~minutes. Their findings, including the omitted citation of~\cite{Remondiere2026}, were
  applied before deposit.
\item \textbf{Pre-release audit.} 25 subagents (four auditors for privacy, reproducibility,
  metadata and claims, and 21 adversarial verifiers): about $4.2\times10^6$ tokens, 473 tool calls,
  about 19~minutes. It ran on 10 November 2026, after the release, with web access. It rebuilt the
  Lean project and the paper from a clean export of the then-current commit, and found that OpenAI's
  Lean code had been released, which made this paper's earlier pre-registration claim false. To
  confirm the release, its agents retrieved the released repositories' metadata, file listings and
  README files; none opened OpenAI's Lean files or the committee's assessment. Its confirmed findings
  were applied before deposit. Its result record is not included, because it quotes the private
  details it was searching for.
\item \textbf{Post-release wording check.} 21 subagents (four checkers for facts, consistency,
  honesty and privacy, and 17 adversarial verifiers), instructed not to open OpenAI's Lean code or
  the assessment: about $3.8\times10^6$ tokens, 394 tool calls, about 14~minutes. Its confirmed
  findings were applied before deposit.
\end{itemize}
The orchestration scripts, which contain every subagent prompt, are included in the repository
(\lean{dev/workflows/}).

\paragraph{Verification and its limits.}
\begin{itemize}[leftmargin=*]
\item Every result in Table~\ref{tab:checked} is checked by Lean's kernel when the project is
  built, and \lean{Audit.lean} reproduces the axiom report. No independent replay (for example
  with \lean{lean4checker} or Comparator) has yet been run.
\item Every bibliography entry was checked by an AI verifier for existence and bibliographic
  metadata against arXiv, Crossref or the cited URL. This does not mean that the cited content was
  read. Entries that could not be confirmed are flagged.
\item Natural-language claims were reviewed only by AI agents, all instances of the same model
  family, whose judgements may be correlated; errors may remain.
\end{itemize}

\paragraph{Authorship, competing interests and funding.}
AI systems are not listed as authors, in line with prevailing publication norms; responsibility
rests with the human author. The author declares no employment by, funding from, or financial
interest in OpenAI, Anthropic or any other AI developer, and received no funding for this work. Use
of the AI systems was paid for by the author through a commercial subscription. Anthropic, whose
models were used, had no role in the design, conduct or reporting of this work.

%=====================================================================
\section{Data and code availability}\label{sec:data}
%=====================================================================
The complete package is at \url{https://github.com/betterfuture2030/yang-mills-lean} and in its
Zenodo archive, which is built from the tagged Git commit:
\begin{itemize}[leftmargin=*]
\item the Lean project (\lean{lean/YangMills});
\item the problem description (\lean{PROBLEM.md});
\item the eleven library notes, the completeness critique and the bibliography
  (\lean{library/});
\item the plan and its five critiques (\lean{plan/});
\item the workflow scripts and their run summaries (\lean{dev/workflows/});
\item this paper (\lean{paper/}).
\end{itemize}
Third-party PDFs are not redistributed. A manifest with arXiv identifiers and a download script
(\lean{scripts/fetch\_pdfs.sh}) is provided instead.

\paragraph{Statement-file hashes.} SHA-256 of the statement files of version~0.2.0:
\begin{Verbatim}
lean/YangMills/YangMills/OS.lean
  36b5f3695467ed0206d12a70c0705b669d34cf9df9695982ef521eecfd546a33
lean/YangMills/YangMills/Lattice.lean
  f5807c28379cfcea5175e0fc14205b590792f89e3f7b10613cf6b77e1c1637ce
lean/YangMills/YangMills/SU.lean
  876d8e2c064c30c71a890bd243b1e5df4e91c994b7a1d4824eb3bfebe7dc8c1e
lean/YangMills/YangMills/Statement.lean
  9cbfec13626e1c3eab5f68d91a598f50765f6c3b9a3d45b5ecece2dc29a5bd05
lean/YangMills/YangMills/OpenProblems.lean
  99751a4eb989696b4c17c1eb4148643480f5786451e338df3b71f893746c009b
\end{Verbatim}

To reproduce the audit:
\begin{Verbatim}
cd lean/YangMills && lake exe cache get && lake build
lake env lean YangMills/Audit.lean
\end{Verbatim}
Text is licensed under CC~BY~4.0 and Lean code under Apache-2.0.

\section*{Acknowledgements}
The author thanks the authors and maintainers of Lean, Mathlib and the public formalizations and
manuscripts cited here; none of them was consulted about, or is responsible for, this work.

\bibliographystyle{unsrtnat}
\bibliography{refs}

\end{document}
